% !TeX program = lualatex
% !TeX encoding = utf8
% !TeX spellcheck = uk_UA
% !TeX root = ../ElMagConspect.tex

%=========================================================
\chapter{Електричні заряди та їх взаємодія}\label{\currfilebase}
\graphicspath{{\currfilebase/Pictures}}
%=========================================================
\epigraph{\Alexander  ...взаємне тяжіння електричної рідини, іменованої позитивною, і електричної рідини, іменованої зазвичай негативною,
	полягає у зворотному відношенні квадратів відстаней...
}{Шарль Огюстен Кулон}


%% --------------------------------------------------------
\section{Сили (взаємодії) в природі}
%% --------------------------------------------------------

Фізичні явища в природі зумовлені взаємодією між тілами. Наприклад, рух Землі навколо Сонця пояснюється гравітаційною взаємодією, яка визначає орбіту
планети. Магнітні сили притягують залізні предмети, а молекулярні сили дозволяють воді підніматися в капілярі. $\beta$-Радіоактивний розпад
ядер є результатом слабкої взаємодії, тоді як притягання між нуклонами в ядрі відбувається завдяки сильній ядерній взаємодії. Незважаючи на
різноманіття
цих явищ, усі взаємодії в природі можна звести до чотирьох фундаментальних типів: гравітаційної, електромагнітної, сильної (ядерної) та слабкої
взаємодій.


Гравітаційна та електромагнітна взаємодії досліджуються в рамках класичної механіки, хоча для електромагнітних сил, на відміну від гравітаційних,
розроблена квантова теорія, відома як квантова електродинаміка. Водночас сильна та слабка взаємодії описуються виключно в
рамках квантової теорії. Ці чотири типи взаємодій охоплюють весь спектр фізичних процесів --- від подій у Всесвіті до явищ на рівні елементарних
частинок та процесів у живих організмах (рис.~\ref{pic:interactions}).

\begin{figure}[h!]\centering
	\includegraphics[width=\linewidth]{forces}
	\caption{Взаємодії в природі}
	\label{pic:interactions}
\end{figure}


%Основні типи фізичних полів:
%\begin{enumerate}
%	\item Гравітаційне поле забезпечує взаємодію між масами.
%	\item Електромагнітне поле діє між зарядами та струмами.
%	\item Глюонне поле відповідає за сильну взаємодію.
%	\item Поле векторних бозонів опосередковують слабку взаємодію.
%\end{enumerate}

Серед чотирьох фундаментальних взаємодій електромагнітна займає особливе місце завдяки своєму впливу на всі рівні організації матерії --- від атомів до
макроскопічних об’єктів.\footnote{Сьогодні існують припущення щодо існування матерії, яка не підкоряється електромагнітній взаємодії і тому не випромінює та не поглинає світло.
Така матерія називається темною. Її існування можна виявити через гравітаційний вплив на звичайну матерію.} Вона визначає структуру атомів і молекул, а
також властивості речовини, які ми спостерігаємо у повсякденному житті. Саме
завдяки електромагнітній взаємодії атоми утримуються разом, забезпечуючи існування твердих і рідких тіл.

Електромагнітна взаємодія також є основою для таких явищ, як поширення світла, радіохвиль та інших форм випромінювання. Ці процеси відіграють ключову
роль у передачі енергії та інформації, що є невід’ємною частиною сучасного життя.

Звичайні сили, які ми відчуваємо у повсякденному житті, такі як тиск руки на камінь чи взаємодія підошви із землею, також обумовлені електромагнітною
взаємодією. Коли два тіла стикаються, їхні атоми зближуються на відстані порядку атомних розмірів, і взаємодія між електронами та ядрами цих атомів
визначає характер цього контакту, створюючи силу, яка не дозволяє тілам проникати одне в одне.

Особливістю електромагнітної взаємодії є її величезна сила порівняно з гравітаційною. Навіть невеликий дисбаланс між електричними зарядами може
призвести до ефектів, що перевершують найпотужніші гравітаційні явища. Як яскраво зазначає Річард Фейнман у своїх лекціях:
\begin{tornpage}\small\itshape%\Alexander
	\hspace{\parindent}Якби у вашому тілі або в тілі вашого сусіда (який стоїть від вас на відстані витягнутої руки) електронів виявилося б всього на 1\%
	більше, ніж протонів, то сила вашого відштовхування була б неймовірно великою. Наскільки великою? Достатньою, щоб підняти хмарочос? Більше!
	Достатньою,
	щоб підняти гору Еверест? Більше! Сили відштовхування вистачило б, щоб підняти «вагу», що дорівнює вазі нашої Землі!\\

	\hfill{ Р. Фейнман у книзі <<Фейнманівські лекції з фізики>>}.
\end{tornpage}
%В основі всіх фізичних явищ лежить взаємодія між тілами або частинками, що беруть участь у цих явищах. Наприклад, Земля рухається навколо Сонця через
%гравітаційну взаємодію з ним. Іншим прикладом взаємодії є притягання або відштовхування електричних зарядів, яке є проявом електричної взаємодії.
%Відштовхування або притягання магнітних полюсів чи струмів --- це прояв магнітної взаємодії. У XIX столітті Джеймс Клерк Максвелл у своїй теорії
%об'єднав електричну та магнітну взаємодію. Це об'єднання отримало назву електромагнітної взаємодії.
%
%У мікросвіті, крім електромагнітної та гравітаційної взаємодій, діють ще два типи: слабка та сильна ядерні взаємодії. Сильна взаємодія утримує кварки
%всередині протонів і нейтронів, а також протони й нейтрони у ядрах атомів. Слабка взаємодія відповідає за ядерні перетворення, такі як бета-розпад.


%Всяка взаємодія передається через поле. Щоб зрозуміти це, слід звернути увагу на поняття <<далекодії>> та <<близькодії>>. У випадку далекодії взаємодія
%між тілами відбувається миттєво на будь-якій відстані, без явного посередника. Це уявлення панувало до XIX століття. Однак із розвитком фізики стало
%зрозуміло, що жодна взаємодія не може передаватися без проміжного агента. Таким агентом є поле. Земля взаємодіє із Сонцем через гравітаційне поле.
%Електричні заряди та струми взаємодіють через електромагнітне поле. Кварки в середині протона чи нейтрона взаємодіють через глюонне поле. А розпад
%протона, чи взаємодія електрона та нейтрино відбувається завдяки
%$W^{\pm}$- та $Z$-бозонним полям. Кожна частинка створює поле, і це поле вже діє на іншу частинку.
%
%
%Взаємодії у природі характеризуються кількома важливими властивостями: їхнім джерелом, ефектом, радіусом дії та інтенсивністю. Розглянемо ці аспекти
%детальніше, щоб зрозуміти основні відмінності між чотирма фундаментальними типами взаємодій.
%
%\emphx{Джерело взаємодії.} Кожна взаємодія виникає завдяки певним властивостям матерії. Наприклад, гравітаційна взаємодія обумовлена масою,
%електромагнітна --- електричним зарядом, слабка ядерна --- слабким зарядом, а сильна ядерна --- кольоровим зарядом кварків.
%
%\emphx{Ефект взаємодії.} Ця характеристика описує, як взаємодія впливає на тіла або частинки. Гравітація змушує всі тіла притягуватися одне до одного.
%Електромагнітна взаємодія може викликати як притягання, так і відштовхування залежно від знаків зарядів. Слабка ядерна взаємодія спричиняє ядерні
%перетворення, тоді як сильна взаємодія утримує кварки разом у складі протонів і нейтронів.
%
%\emphx{Радіус дії.} Ця властивість визначає, на якій відстані взаємодія залишається відчутною. Гравітація та електромагнітна взаємодія діють на
%необмежених відстанях, хоча їхній вплив слабшає із зростанням відстані. Слабка та сильна ядерні взаємодії мають дуже обмежений радіус дії, що становить
%$10^{-16}$ см і $10^{-13}$ см відповідно.
%
%\emphx{Інтенсивність.} Інтенсивність показує силу взаємодії. Гравітація є найслабшою з усіх типів взаємодій, тоді як сильна ядерна взаємодія є
%найпотужнішою. Електромагнітна взаємодія займає проміжне місце між цими крайнощами, а слабка ядерна взаємодія за своєю інтенсивністю перевищує
%гравітацію, але значно поступається електромагнітній і сильній.
%
%Для зручності всі ці характеристики зібрано в таблицю~\ref{tab:forces}.



%\emphx{Електромагнітна взаємодія} є одним із фундаментальних типів
%взаємодій і передається за допомогою агента ---
%\emphx{електромагнітного поля}.
%
%\emphx{Електромагнітним полем} називається область простору, де діють
%електричні та магнітні сили.

%Поле \emphx{створюється} електричними \textit{зарядами} і \emphx{діє}
%на \textit{заряди}.


\section{Електричний заряд}

\emphx{Електричний заряд} --- це фундаментальна фізична величина, яка визначає інтенсивність взаємодії між зарядженими тілами. Він є властивістю деяких
елементарних частинок, завдяки якій ці частинки взаємодіють через електромагнітну взаємодію. Фундаментальність заряду означає, що це базова
характеристика частинок, яка не залежить від інших властивостей і є невід'ємною складовою їхньої природи. Оскільки заряд є такою основною величиною, ми не
можемо вивести його властивості з інших, більш фундаментальних принципів. Натомість ми змушені описувати його властивості через спостереження та
експериментальні дослідження, що дозволяє нам виявити характер і поведінку заряду в різних умовах. Зокрема, експериментально було виявлено наступні
властивості:

\begin{enumerate}
	\item \emphx{Дискретність заряду}.

    Експеримент, проведений у 1909 році Робертом Ендрюсом Міллікеном та Гарві Флетчером, дозволив встановити, що заряд має дискретну природу. В
    результаті досліджень було визначено величину заряду електрона з точністю до 1\%, що підтвердило існування мінімальної одиничної величини заряду.

    Згідно з цими результатами, будь-який ненульовий заряд є кратним деякому елементарному заряду, числове значення якого дорівнює заряду електрона:

    \begin{equation*} e = 4.803 \cdot 10^{-10}\ \text{Фр} = 1.602 \cdot 10^{-19}\ \text{Кл}. \end{equation*}

    Це означає, що заряд не може набувати довільних значень, а лише кратних елементарному заряду.

        Крім того, в дослідах Франкліна було встановлено, що в природі існують \emphx{заряди двох типів}, які умовно ділять на \emphx{позитивні} та
	      \emphx{негативні}. Позитивний заряд несуть, наприклад, протони, а негативний --- електрони.

%    \item Адитивність заряду. Заряд є адитивною величиною. Це означає

	\item \emphx{Збереження заряду}.

	      Заряд --- це величина, що \emphx{зберігається}. Позитивні й негативні заряди можуть народжуватися в рівних кількостях або взаємно
	      знищувати один одного. При цьому сумарний електричний заряд замкненої системи не змінюється.
	\item \emphx{Релятивістська інваріантність заряду}.

	      Величина електричного заряду однакова відносно різних систем відліку. Іншими словами, величина заряду не залежить від швидкості.
\end{enumerate}


%% --------------------------------------------------------
\section{Як заряджають тіла?}
%% --------------------------------------------------------


Тіла складаються з заряджених частинок, і в нормальних умовах їхня кількість є збалансованою, що робить тіло електрично нейтральним. Однак, коли два
тіла вступають у контакт, між ними може відбутися перенесення електронів, що змінює баланс зарядів. Цей процес інтенсифікується при терті між тілами, що
призводить до ще більшого перенесення електронів. Таке явище називається \emphx{трибоелектричним ефектом} (рис.~\ref{pic:triboeffect}). В результаті
цього ефекту одне тіло набуває негативного заряду, а інше --- позитивного.

%---------------------------------------------------------
\begin{SCfigure}[][h!]
	\centering
	
\includegraphics[height=3cm]{charging-by-friction}
	\caption{Демонстрація трибоелектричного ефекту.\label{pic:triboeffect}}
\end{SCfigure}
%---------------------------------------------------------

Слово <<\emphx{електрика}>> походить від грецького
<<$\eta\lambda\varepsilon\kappa\tau\rho o\nu$>>
(\emphx{електрон}), що означає <<\emphx{бурштин}>>. Люди помітили,
що бурштин, натертий об хутро, притягує дрібні предмети, і це
явище згодом стало основою для дослідження електричних явищ.



Матеріали, що виявляють трибоелектричний ефект, прийнято
розташовувати в \emphx{трибоелектричний ряд}, один кінець якого є
позитивним, а інший --- негативним. Під час тертя пари матеріалів
з ряду, матеріал, розташований ближче до позитивного кінця ряду,
зарядиться позитивно, а інший --- негативно (рис.~\ref{tikz:triboseries}).


\begin{figure}[h!]
	\centering
	\begin{tikzpicture}
		% Градієнт фону
		\shade[left color=blue!30, right color=red!30] (-5.5, -.5) rectangle (5.5, -3.5);

		\node[single  arrow, single  arrow head extend=.5cm, left color=blue!50, right color=red!50, text=white, font={\bf \large}, inner sep=0.3cm]
		(doublearrow)
		{--\hspace*{10.5cm}+};

		% Трибоелектричний ряд
		\node[align=center, rotate=90] at (5,  -2) {Хутро};
		\node[align=center, rotate=90] at (4,  -2) {Фланель};
		\node[align=center, rotate=90, align=center, text width=2.5cm] at (3,  -2) {Слонова\\ кістка};
		\node[align=center, rotate=90] at (2,  -2) {Пір'я};
		\node[align=center, rotate=90] at (1,  -2) {Флінтглас};
		\node[align=center, rotate=90] at (0,  -2) {Шовк};
		\node[align=center, rotate=90] at (-1, -2) {Алюміній};
		\node[align=center, rotate=90] at (-2, -2) {Поліестер};
		\node[align=center, rotate=90] at (-3, -2) {Папір};
		\node[align=center, rotate=90] at (-4, -2) {Бурштин};
		\node[align=center, rotate=90] at (-5, -2) {Ебоніт};
	\end{tikzpicture}
	\caption{Трибоелектричний ряд\label{tikz:triboseries}}
\end{figure}


Трибоелектричний ефект має не лише пізнавальне, а й практичне значення. З одного боку, він лежить в основі \emphx{трибоелектричних наногенераторів} (TENG) --- пристроїв, які перетворюють механічну енергію тертя чи вібрацій (кроки людини, вітер, хвилі) в електричну; завдяки високій ефективності при низьких частотах їх розглядають як перспективне джерело живлення для носимої електроніки та автономних сенсорів у системах <<інтернету речей>>.

З іншого боку, неконтрольоване трибоелектричне заряджання створює практичні проблеми. У мікроелектроніці раптовий електростатичний розряд (\emphx{ESD}) здатний вивести з ладу чутливі напівпровідникові елементи, тому виробництво мікросхем супроводжується суворими заходами захисту від статичної електрики (антистатичні браслети, заземлені робочі поверхні, іонізатори повітря). У хімічній та харчовій промисловості трибоелектричне заряджання дрібнодисперсних матеріалів (борошна, вугільного пилу) може спричинити електростатичний розряд --- одну з причин вибухів пилоповітряних сумішей. Водночас керований трибоелектричний ефект використовують і на користь: в \emphx{електрофільтрах} заряджені частинки диму чи пилу осідають на протилежно заряджених електродах, очищуючи промислові гази, а в \emphx{ксерографії} (принцип роботи копіювальних апаратів і лазерних принтерів) електростатичне заряджання світлочутливого барабана й тонера забезпечує перенесення зображення на папір.


%% --------------------------------------------------------
\section{Закон Кулона}
%% --------------------------------------------------------


Експериментально встановлено, що заряди одного знака (однойменні заряди) \emphx{відштовхуються}, а заряди різного знака (різнойменні заряди) ---
\emphx{притягуються}. Ці властивості зарядів лягли в основу формулювання закону їхньої взаємодії.


Для спрощення аналізу часто використовують ідеалізовані моделі заряджених тіл, однією з яких є \emphx{точковий заряд}. Точковим називається такий заряд,
розмірами якого в конкретних умовах можна знехтувати.

Із цією моделлю заряджених тіл пов'язаний важливий закон, встановлений Ш.~Кулоном у 1785 році. Він стверджує, що \emphx{сила, що діє між двома
нерухомими точковими зарядами $q_1$ і $q_2$, розташованими на відстані $r$ один від одного, описується законом}:
\begin{equation}\label{eq:ColumbLaw}
	\tcbhighmath{
		\vect{F}_{12} = \frac{q_1 q_2}{r^2}\frac{\vect{r}_{12}}{r}.
	}
\end{equation}
Ця сила напрямлена вздовж лінії, яка з'єднує заряди, тобто носить \emphx{центральний} характер.
\begin{Attention}
    У реальних умовах точкові заряди не існують, і для застосування закону Кулона до малих заряджених тіл необхідно, щоб заряд всередині них був
    розподілений рівномірно, інакше порушується центральний характер взаємодії. Тому під <<фізично>> точковими зарядами будемо розуміти такі, в яких
    розподіл заряду
    є рівномірним,
    тобто не залежить від орієнтації тіла.
\end{Attention}

На рис.~\ref{pic:Columb} показані величини, що входять до закону Кулона, підкреслено центральний характер взаємодії та її напрям однойменних та
різнойменних зарядів.


\begin{SCfigure}[0.5][h!]
	\centering
	\localinput{Columb.tikz}
	\caption{Ілюстрація до закону Кулона}
    \label{pic:Columb}
\end{SCfigure}

Досліди самого Кулона не відзначалися високою точністю, і можна було припустити, що знайдений ним простий закон обернених квадратів виконується лише
наближено, наприклад:
\begin{equation*}
	F = \frac{q_1 q_2}{r^{2 + \epsilon}},
\end{equation*}
де $|\epsilon| \ll 1$.

У випадку, коли $\epsilon = 0$, заряди зосереджуються лише на поверхні провідника, а всередині провідника їх не існує. Якщо припустити $|\epsilon| \neq
0$, то це може спричинити появу заряду в товщі провідника.

\begin{SCfigure}[0.75][h!]
	\centering
	\includegraphics[width=0.5\linewidth]{cavindesh_exp}
	\caption{Установка Г.~Кавендіша для перевірки закону обернених квадратів.}
	\label{pic:Cavendih_exp}
\end{SCfigure}

Експеримент, що дозволив перевірити ці припущення, був проведений Г.~Кавендішем і полягав у накладанні двох провідних напівсфер на заряджену провідну
кулю, що щільно прилягали одна до одної, утворюючи майже суцільну оболонку (рис.~\ref{pic:Cavendih_exp}). За умов закону обернених квадратів весь заряд
мав би перейти на зовнішні напівсфери. Якщо ж взаємодія описується за формулою з $\epsilon \neq 0$, то всередині кулі залишався б
ненульовий заряд, який пропорційно зменшується зі зменшенням значення $\epsilon$. Цей заряд можна було б виміряти за допомогою чутливого електрометра,
приєднаного до внутрішньої сфери (на рисунку роль електрометра відіграють дві кульки на нитках). Експеримент повторювався з поступовим підвищенням
точності вимірювань. У 1983 році було отримано обмеження для
$\epsilon$:
\begin{equation*}
	|\epsilon| < 10^{-16} \text{–} 10^{-17}.
\end{equation*}

\begin{Attention}
Порушення закону обернених квадратів мало б серйозні наслідки для теорії електромагнітної взаємодії, оскільки цей закон є основою для її опису. Теорія
Максвелла, яка визначає поведінку електричних та магнітних полів, спирається на його точність. Якби закон порушувався, це змінило б наші уявлення про
електромагнітну взаємодію та вимагало б перегляду численних експериментальних результатів, що сьогодні підтверджують цю теорію.
\end{Attention}

%Закон Кулона знаходить застосування далеко за межами класичної електростатики. У \emphx{фізиці твердого тіла} він лежить в основі розрахунку електростатичної енергії кристалічної ґратки: підсумовування кулонівських взаємодій між усіма іонами кристала дає \emphx{сталу Маделунга}, яка визначає внесок електростатики в енергію зв'язку іонних кристалів (наприклад, NaCl) і, відповідно, впливає на такі властивості матеріалів, як температура плавлення та механічна міцність.
%
%У \emphx{мікро- та наноелектромеханічних системах} (МЕМС) кулонівська взаємодія між зарядженими електродами використовується для керованого переміщення мікроскопічних елементів --- на цьому принципі побудовані електростатичні мікроактюатори, що застосовуються, зокрема, в акселерометрах смартфонів, мікродзеркальних матрицях проєкторів (DMD) та мікронасосах. Подібний принцип лежить і в основі \emphx{іонних пасток} --- пристроїв, які утримують заряджені частинки в обмеженому об'ємі за допомогою електричних (та магнітних) полів; іонні пастки є ключовим інструментом сучасної метрології (атомні годинники) і однією з платформ для побудови квантових комп'ютерів.


%% --------------------------------------------------------
\begin{Questions}
\begin{quizlist}
	\item Назвіть чотири фундаментальні взаємодії. Яка з них визначає будову атомів і молекул та силу, з якою рука тисне на камінь? Чому цю силу вважають електромагнітною за природою?
	\item Чому в повсякденному житті ми не відчуваємо величезних електричних сил, про які говорить Фейнман? Що забезпечує електричну нейтральність макроскопічних тіл?
	\item Що означає, що електричний заряд є фундаментальною величиною? Чому його властивості ми вивчаємо з досліду, а не виводимо з інших принципів?
	\item Сформулюйте три основні властивості заряду: дискретність, збереження та релятивістську інваріантність. Чи може тіло мати заряд $2{,}5\,e$? Відповідь обґрунтуйте.
	\item У чому полягав дослід Міллікена та Флетчера і що він довів? Який зв'язок між значеннями елементарного заряду в одиницях СГС (Фр) та СІ (Кл)?
	\item Що таке трибоелектричний ефект і чому при терті двох тіл одне заряджається позитивно, а інше --- негативно? Користуючись трибоелектричним рядом (рис.~\ref{tikz:triboseries}), визначте знаки зарядів, яких набудуть шовк і поліестер при терті один об одного.
	\item Наведіть по одному прикладу корисного та шкідливого прояву трибоелектричного ефекту. Поясніть, як працює електрофільтр або ксерографія.
	\item Сформулюйте закон Кулона й запишіть його у векторній формі. Що означає, що кулонівська сила є \emphx{центральною}?
	\item За яких умов заряджене тіло можна вважати точковим? Чому для «фізично» точкових зарядів потрібен рівномірний розподіл заряду?
	\item Як зміниться сила взаємодії двох точкових зарядів, якщо: а)~відстань між ними збільшити втричі; б)~один із зарядів збільшити вдвічі; в)~обидва заряди змінити знак на протилежний?
	\item У чому полягала ідея досліду Кавендіша? Чому за закону обернених квадратів ($\epsilon = 0$) заряд зосереджується лише на зовнішній поверхні провідника, а при $\epsilon \neq 0$ з'явився б заряд усередині?
	\item Яке обмеження на $|\epsilon|$ отримано в експериментах і чому порушення закону обернених квадратів було б настільки важливим для електродинаміки?
\end{quizlist}
\end{Questions}